\chapter{Calculators \& File I/O}
\label{ch:calculators}

\newthought{Calculators are the external programs} --- spectrum generators, event
generators, analysis tools --- that \Jtwo\ exists to orchestrate. Each entry under
\yamlkey{Calculators.Modules} declares one program: where it lives, how to isolate it,
what environment it needs, and the \emph{write input $\rightarrow$ run commands
$\rightarrow$ read output} cycle every Sample drives it through, including the file
formats that cycle reads and writes. This is the longest block in a real card and the
longest chapter in this part.

\section{The execution model}
\label{sec:calc-model}

Each Worker holds \textbf{one instance per calculator type} for its whole life:
templates are parsed once, \yamlkey{env\_setup} scripts are sourced once, and every
Sample the Worker handles reuses them. Per run, the cycle is fixed:

\begin{enumerate}
  \item \textbf{Acquire a slot.} The Worker blocks on the calculator's global
        concurrency pool and receives a \emph{pack id}
        (\S\ref{sec:pools},~\S\ref{sec:clone-shadow}).
  \item \textbf{Prepare the runtime.} Shadow install (first time a Worker sees this
        pack) or a symlink into the Sample directory.
  \item \textbf{Write inputs} --- each \yamlkey{execution.input} spec, through \Portal\
        (Section~\ref{sec:input-spec}).
  \item \textbf{Run commands} --- each \yamlkey{execution.commands} entry, sequentially,
        under the module's timeout.
  \item \textbf{Read outputs} --- each \yamlkey{execution.output} spec; captured
        variables merge into the observables.
  \item \textbf{Release the slot} --- always, success or failure.
\end{enumerate}

A non-zero exit code, a timeout, or an I/O error fails the Sample (never the scan) and
the slot is released.

\section{Global concurrency: \yamlkey{Calculators.Pools}}
\label{sec:pools}

\begin{yamlwin}
Calculators:
  Pools:            # alias: pools
    SPheno: 8       # at most 8 concurrent SPheno runs, across ALL Workers
    MadGraph: 2
\end{yamlwin}

Pools are global caps enforced through Redis: a Worker needing a slot blocks until one
frees anywhere in the system (up to an internal 30\,s acquire timeout, after which the
Sample fails). Values have floor 1. When a module has no pool entry, its per-module
\yamlkey{make\_parallel} value (default 1) is used --- itself a legacy knob \cli{Jarvis
man} now points away from, in favour of \yamlkey{Pools}.

\begin{jtip}
Give every expensive calculator an explicit pool sized to what the machine actually
sustains --- licenses, memory, scratch bandwidth. Pools are the throttle that lets you
run more Workers than any single tool could survive.
\end{jtip}

\section{\yamlkey{Modules[]}: the per-module keys}
\label{sec:calc-module-keys}

Every entry in the list is described by the keys of Table~\ref{tab:calc-module-keys}.

\begin{table}[ht]
  \footnotesize
  \begingroup
  \setlength{\tabcolsep}{0pt}
  \begin{tabular}{@{}p{0.188\linewidth}p{0.104\linewidth}p{0.188\linewidth}p{0.52\linewidth}@{}}
    \toprule
    \textbf{Key} & \textbf{Req.} & \textbf{Default} & \textbf{Meaning} \\
    \midrule
    \yamlkey{name}               & \req & ---     & module name; nameless entries are silently dropped \\
    \yamlkey{required\_modules}  &      & \code{[]} & dependencies $\rightarrow$ layers (Ch.~\ref{ch:workflow-model}) \\
    \yamlkey{clone\_shadow}      &      & \code{false} & physical per-pack copy for non-thread-safe tools \\
    \yamlkey{path}               &      & \yamlkey{execution.path} $\to$ \code{"."} & runtime directory; may contain \token{@PackID} when shadowed \\
    \yamlkey{source}             &      & ---     & shadow: copy source; non-shadow: tool symlinked into the Sample dir \\
    \yamlkey{symlink\_name}      &      & module name & name of that symlink \\
    \yamlkey{timeout}            &      & none    & seconds; one deadline across \emph{all} execution commands \\
    \yamlkey{make\_parallel}     &      & \code{1} & legacy pool fallback when no \yamlkey{Pools} entry exists (\S\ref{sec:pools}) \\
    \yamlkey{env\_setup}         &      & \code{[]} & environment scripts (\S\ref{sec:env-setup}) \\
    \yamlkey{installation}       &      & \code{[]} & shadow install commands, once per pack (\S\ref{sec:clone-shadow}) \\
    \yamlkey{initialization}     &      & \code{[]} & post-install commands, once per pack \\
    \yamlkey{execution}          & \req & \code{\{\}} & commands + I/O specs, required in practice (\S\ref{sec:execution-block}) \\
    \bottomrule
  \end{tabular}
  \endgroup
  \caption{\yamlkey{Calculators.Modules[]} keys.}
  \label{tab:calc-module-keys}
\end{table}

\section{Isolation: safe tools vs \yamlkey{clone\_shadow}}
\label{sec:clone-shadow}

Two isolation strategies, chosen per module:

\paragraph{Concurrency-safe tools} (default). The installed tool is shared; each Sample
gets a \emph{symlink} to it inside the Sample directory (named \yamlkey{symlink\_name}),
so commands can call it with a stable relative name and nothing is copied. Combine with
\yamlkey{LibDeps.registered\_executables} (Chapter~\ref{ch:libdeps}) to skip even the
symlink.

\paragraph{Non-thread-safe tools: \yamlkey{clone\_shadow: true}.} Tools that scribble on
their own installation directory (\textsc{madgraph}-style generators) get \emph{physical
copies}. Each concurrency slot --- each \emph{pack} --- owns a private directory tree:

\begin{yamlwin}
- name: MadGraph
  clone_shadow: true
  source: "&J/External/MG5"          # copied per pack when no installation given
  path: "&J/packs/MG5-@PackID"       # each slot's private tree
  installation:
    - cmd: "cp -r \${source}/* \${path}/"
      cwd: "\${path}"
  initialization:
    - cmd: "./bin/compile"
      cwd: "\${path}"
\end{yamlwin}

Pack ids are \textbf{stable slots} (\code{001}\,\ldots\,\code{N}, one per pool slot)
whose directories persist across Samples, Workers, and runs. The tokens
\code{\$\{source\}} and \code{\$\{path\}} are available inside
\yamlkey{installation} / \yamlkey{initialization} commands only.

\begin{jwarn}
Because packs persist, \textbf{installation commands must be idempotent}: every new
Worker re-runs \yamlkey{installation} on a possibly already-populated pack directory. A
plain overwrite-copy (as above) is fine; an \code{append}, a \code{patch}, or an
unconditional \code{configure} that fails on a second run is not.
\end{jwarn}

\section{Environments: \yamlkey{env\_setup}}
\label{sec:env-setup}

Tools that need a sourced environment (\textsc{rivet}, \textsc{root}-based stacks)
declare it per module:

\begin{yamlwin}
- name: RivetAnalysis
  env_setup:
    - source: "&J/External/Rivet/rivet_env.sh"
  execution:
    commands:
      - cmd: "rivet --analysis=MY_ANALYSIS input.hepmc"
        cwd: "@Sdir"
\end{yamlwin}

At first use, the Worker runs \cli{source <script> \&\& env} in an isolated shell,
captures the resulting environment, merges it over its own, and \textbf{caches it for the
process lifetime} --- each script runs once per Worker, not once per Sample. All
subsequent commands of that module run under the captured environment.

\section{The \yamlkey{execution} block}
\label{sec:execution-block}

\begin{yamlwin}
execution:
  path: "&J/calc/eggbox"       # fallback for the module path
  commands:
    - cmd: python3 eggboxlk    # registered names + all tokens resolve here
      cwd: "@Sdir"
  input:
    - name: inpjson
      path: "@Sdir/input.json"
      type: JSON
      actions:
        - type: Dump
          variables:
            - name: xx
              expression: x * Pi     # write a computed value
            - name: yy               # no expression -> copy observable yy
  output:
    - name: oupjson
      path: "@Sdir/output.json"
      type: JSON
      variables:
        - name: z
          entry: z                   # dotted path into the file
\end{yamlwin}

\begin{keylist}
  \item[commands] A list of \code{\{cmd, cwd\}} maps, run \emph{sequentially} through
        the Worker's scheduler; \yamlkey{cwd} defaults to \code{"."} relative to the
        module path. Registered executable names, static tokens, and per-Sample tokens
        all resolve inside \yamlkey{cmd} and \yamlkey{cwd}.
  \item[input] File specs written \emph{before} the commands run.
  \item[output] File specs read \emph{after} the commands finish.
\end{keylist}

Both lists share one spec format, and that format is the subject of the rest of this
chapter. Notice what a spec does \emph{not} say: nothing about how a \code{JSON} or an
\code{SLHA} file is actually parsed. A spec names a \yamlkey{type} and the reading and
writing are delegated to the \Portal\ package --- which is why upgrading \Portal\ adds
formats without a runtime release.

\section{The division of labour}
\label{sec:io-ownership}

\begin{itemize}
  \item \textbf{\Jtwo\ owns the semantics}: the spec syntax, token resolution in
        \yamlkey{path}, expression evaluation in \yamlkey{Dump} variables, and the
        \yamlkey{save} policy that copies files into \file{SAMPLE/}.
  \item \textbf{\Portal\ owns the formats}: how a \code{JSON} or \code{SLHA} file is
        actually serialized and parsed --- \emph{variable read/write only}, no file
        management.
\end{itemize}

\begin{jwarn}
An unknown or empty \yamlkey{type} is a \textbf{hard failure} that lists the formats
your installed \Portal\ registers --- never a silent skip. If a format you expect is
missing from that list, upgrade \Portal.
\end{jwarn}

\section{The format catalog}
\label{sec:format-catalog}

Table~\ref{tab:portal-formats} lists the \yamlkey{type} values \Portal\ understands
today.

\begin{table}[ht]
  \footnotesize
  \begingroup
  \setlength{\tabcolsep}{0pt}
  \begin{tabular}{@{}p{0.34\linewidth}p{0.66\linewidth}@{}}
    \toprule
    \textbf{type} & \textbf{Format} \\
    \midrule
    \code{JSON}    & \textsc{json} documents; dotted \yamlkey{entry} paths navigate nesting \\
    \code{CSV} / \code{TSV} & delimited tables \\
    \code{DAT}     & whitespace-separated numeric files \\
    \code{Wolfram} & Mathematica expression files \\
    \code{SLHA} / \code{xSLHA} & \textsc{susy} Les Houches Accord (and extended) spectrum files \\
    \bottomrule
  \end{tabular}
  \endgroup
  \caption{Formats exposed by the current \Portal\ V2 surface. Third-party packages can
  register more through the \code{jarvishep.io} entry point; the authoritative list is
  \cli{Jarvis portal man}.}
  \label{tab:portal-formats}
\end{table}

\section{\yamlkey{input[]}: writing calculator input}
\label{sec:input-spec}

\begin{yamlwin}
input:
  - name: lha_in
    path: "@Sdir/LesHouches.in"
    type: SLHA
    save: true                    # keep a copy under SAMPLE/
    actions:
      - type: Dump
        variables:
          - name: m0
            expression: m0            # computed (here: verbatim by expression)
          - name: tanb
            expression: "10 * tb_u"   # any shared-language expression
          - name: sign_mu             # no expression -> copy observable verbatim
\end{yamlwin}

\begin{keylist}
  \item[name] The spec's label (used in logs and generated observable names).
  \item[path] Where to write; per-Sample tokens resolve here, and \token{@Sdir}
        materializes the Sample directory.
  \item[type] The \Portal\ format.
  \item[actions] Today exactly one action type exists: \yamlkey{Dump} writes the listed
        variables into the file. Each entry writes \yamlkey{name} either from an
        \yamlkey{expression} (shared language, Chapter~\ref{ch:expressions}) or, with no
        expression, verbatim from the observable of the same name --- a missing
        observable becomes the literal string \code{MISSING\_VALUE}, which is your cue
        in the produced file.
\end{keylist}

If the target file already exists, \yamlkey{Dump} \emph{merges} into it rather than
replacing it --- useful for layering computed values onto a template your commands
copied into place first.

\section{\yamlkey{output[]}: reading calculator output}
\label{sec:output-spec}

\begin{yamlwin}
output:
  - name: spectrum
    path: "@Sdir/SPheno.spc"
    type: SLHA
    save: true
    variables:
      - name: mh
        entry: "MASS.25"          # dotted path into the parsed structure
      - name: mstop1
        entry: "MASS.1000006"
\end{yamlwin}

Each \yamlkey{variables} entry maps \yamlkey{name} from \yamlkey{entry} --- a dotted
path into the parsed file, defaulting to the name itself. A missing entry captures
\code{None} rather than failing, so downstream expressions see a null value they can be
written to tolerate (or fail on, loudly --- Section~\ref{sec:likelihood-failure}).

\section{The \yamlkey{save} policy}
\label{sec:save-policy}

By default nothing a calculator writes or reads is kept --- working files vanish with
the working directory. The \yamlkey{save} flag on a spec keeps a copy under the Sample's
directory, handled by a dedicated per-Worker file-operation process so Workers never
block on slow storage:

\begin{itemize}
  \item \textbf{input specs}: copied into \file{SAMPLE/<bucket>/<uuid>/} only when
        \yamlkey{save: true}.
  \item \textbf{output specs}: \yamlkey{save: true} copies into the Sample directory
        proper; without it, outputs that must survive the working directory go to the
        Sample's hidden \file{.temp/} area instead.
  \item Saved files keep their basename; a name that would collide across modules is
        disambiguated as \code{basename@Module}.
  \item When a spec is saved, the file's final path is also added to the observables,
        so the database records where each kept artifact lives.
\end{itemize}

\begin{jtip}
Save what you would need to re-analyse a point without re-running it --- the spectrum
file, the event-level summary --- and nothing else. At $10^6$ samples, every saved file
is a million files; buckets keep the filesystem sane
(Section~\ref{sec:sample-directory}), but nothing keeps the disk big.
\end{jtip}

\section{Portal from the command line}
\label{sec:portal-cli}

The \Portal\ registry is inspectable and drivable without a scan:

\begin{terminal}
\begin{jcmd}
Jarvis portal man            # list every registered format
Jarvis portal man slha       # one format's manual
Jarvis portal path/to/io.yaml   # run a standalone Portal IO YAML
\end{jcmd}
\end{terminal}

The same commands exist as the standalone \cli{jportal} tool. Running a small
\textsc{io yaml} against a real output file of your calculator is the fastest way to get
\yamlkey{entry} paths right before wiring them into a scan card.

\section{Failure and diagnostics}
\label{sec:calc-failure}

Everything a calculator run does is logged to the Sample's own log: the resolved
command line, raw stdout/stderr, exit status, and timing. On failure the Sample keeps
its working files, so the standard post-mortem is:

\begin{enumerate}
  \item find the failed \textsc{uuid} in the console or \file{samples.csv};
  \item open \file{SAMPLE/<bucket>/<uuid>/Sample\_running.log};
  \item re-run the recorded command by hand in the recorded \yamlkey{cwd}.
\end{enumerate}

\begin{jnote}
The module \yamlkey{timeout} is a deadline across the \emph{whole} command list of one
run, not per command. On expiry the process group receives \textsc{sigterm}, then
\textsc{sigkill}, the Sample fails, and the slot is released --- a hung external tool
cannot wedge the scan.
\end{jnote}
